% Lab 5 -- Exercise 10  |  Nguyễn Huỳnh Hữu Đức – 2510290  |  Original: Ch.6 4.3 (Complementary MOSFET motor controller)  |  Status: Adjust
% Netlist:
%   VDD vdd 0 5 ; VSS vss 0 -5
%   M2 out g vdd vdd SWP l=1u w=1u    (P-channel, symbol pmos4, source + bulk tại +5 V)
%   M1 out g vss vss SW  l=1u w=1u    (N-channel, symbol nmos4, source + bulk tại -5 V)
%   RM out x 20 ; LM x 0 2m           (động cơ = RM nối tiếp LM, giữa out và đất)
%   VIN g 0 PULSE(-5 5 0 1u 1u 1m 2m)
\exercise{Bidirectional Motor Driver with N- and P-channel MOSFETs}{Adjust}{Ch.6 4.3 -- Complementary MOSFET motor controller (optional)}
\statusreason{The original exercise is optional and uses a motor symbol with $\pm V_{DD}$ supplies. It becomes a required
exercise with $\pm 5$\,V supplies (so that $|V_{GS}| = 10$\,V stays within the usual gate rating) and a motor modelled
as $R_M$ in series with $L_M$. Level-1 models replace the PSpice parts. The objective (direction control by high-side
P-channel and low-side N-channel switching) is kept, and the shoot-through current during transitions is added as an observation.}

\task{A P-channel MOSFET connects the motor to $+5$\,V and an N-channel MOSFET connects it to $-5$\,V; both gates are driven by
the same input that switches between $-5$\,V and $+5$\,V. Explain which transistor conducts for each input level,
compute the motor current in both directions, and observe the current drawn from the supplies during the transitions.}

\subsubsection*{Circuit}
\circuitfig{ex10}{Circuit of Exercise 5.10: bidirectional motor driver}
\components{
\begin{itemize}[leftmargin=*,itemsep=1pt]
  \item $V_{DD} = +5$\,V, $V_{SS} = -5$\,V; motor: $R_M = 20\,\Omega$ in series with $L_M = 2$\,mH, between the common drain and ground.
  \item $V_{IN}$: \texttt{PULSE(-5 5 0 1u 1u 1m 2m)} applied to both gates.
  \item M1 (N, symbol \texttt{nmos4}): \texttt{.model SW NMOS(Level=1 Vto=2 Kp=0.5 Lambda=0)}\\
        M2 (P, symbol \texttt{pmos4}): \texttt{.model SWP PMOS(Level=1 Vto=-2 Kp=0.5 Lambda=0)}; $L = W = 1\,\mu$m for both.
\end{itemize}}

\subsubsection*{Theoretical Calculation}
\begin{enumerate}[itemsep=2pt]
  \item For $V_{IN} = -5$\,V and $V_{IN} = +5$\,V, compute $V_{GS}$ of each transistor and state which one is ON.
  \item Compute $R_{DS(\mathrm{on})}$ for $|V_{GS}| = 10$\,V and the steady motor current $I_M$ in each direction.
  \item Compute the time constant $\tau = L_M/(R_M + R_{DS(\mathrm{on})})$.
  \item For which range of $V_{IN}$ are both transistors ON at the same time? What does this mean for the supply current?
\end{enumerate}
\answerspace{12}
\begin{solution}
$V_{IN} = -5$\,V: $V_{GS,P} = -5-5 = -10$\,V (ON), $V_{GS,N} = -5-(-5) = 0$ (OFF): the motor is connected to $+5$\,V.
$V_{IN} = +5$\,V: $V_{GS,P} = 0$ (OFF), $V_{GS,N} = +10$\,V (ON): the motor is connected to $-5$\,V.
\[
  R_{DS(\mathrm{on})} \approx \frac{1}{0.5\cdot(10-2)} = 0.25\,\Omega,\qquad
  I_M = \pm\frac{5}{20+0.25} \approx \pm 247\,\mathrm{mA},\qquad
  \tau = \frac{2\,\mathrm{mH}}{20.25\,\Omega} \approx 99\,\mu\mathrm{s}
\]
Both transistors conduct when $V_{GS,N} > 2$\,V and $V_{GS,P} < -2$\,V, i.e.\ $-3\,\mathrm{V} < V_{IN} < 3\,\mathrm{V}$.
During each 1\,$\mu$s edge the input crosses this range and a large current flows directly from $+5$\,V to $-5$\,V
(shoot-through). Real drivers add a dead time to avoid it.
\end{solution}

\subsubsection*{Simulation}
\begin{enumerate}[itemsep=2pt]
  \item Run the transient analysis. Plot \texttt{V(g)}, \texttt{V(out)} and \texttt{I(RM)}; read $I_M$ in both directions.
  \item Plot \texttt{I(VDD)} and zoom on an input edge to see the shoot-through peak.
\end{enumerate}
\begin{spice}
.model SW  NMOS(Level=1 Vto=2  Kp=0.5 Lambda=0)
.model SWP PMOS(Level=1 Vto=-2 Kp=0.5 Lambda=0)
.tran 0 4m 0 0.1u
\end{spice}
\simfig{ex10_sim}{Input, motor voltage and motor current of Exercise 5.10}

\subsubsection*{Theory vs. Simulation}
\begin{comparison}{Theory vs. simulation, Exercise 5.10}
\cmprow{$I_M$ ($V_{IN}=-5$\,V)}{mA}{$+247$}{$\approx +247$}{$\approx 0$}
\cmprow{$I_M$ ($V_{IN}=+5$\,V)}{mA}{$-247$}{$\approx -247$}{$\approx 0$}
\cmprow{$\tau$}{$\mu$s}{99}{$\approx 100$}{$\approx 1$}
\end{comparison}
\answerspace{3}
\begin{solution}
The steady currents should agree within rounding. The shoot-through peak has no simple hand value; students should
report it from the plot and explain its origin. The motor current settles at about $\pm247$\,mA with a time constant of about
0.1\,ms, and $V(\mathrm{out})$ switches between $-5$ and $+5$\,V, as expected.
\todo{Plot \texttt{Id(M1)} and \texttt{Id(M2)} around $t = 1$\,ms and report the shoot-through peak.}
\end{solution}
